\subsection{Výpočetní náročnost (GPU / CPU / TPU)}
\label{subsec:llm-computational-cost}

Inference jazykového modelu je výpočetně náročná operace.
Její jádro tvoří opakované násobení matic velkých rozměrů v každé vrstvě architektury Transformer, přičemž tyto operace jsou pro každý generovaný token prováděny znovu.
Způsob, jakým je výpočet prováděn, ať už na CPU, GPU nebo specializovaném akcelerátoru jako TPU (Tensor Processing Unit), zásadně ovlivňuje latenci i propustnost systému.
TPU jsou čipy navržené společností Google specificky pro operace s tenzory a mohou nabídnout vyšší efektivitu inference při nižší spotřebě energie, nicméně jejich dostupnost je v současnosti omezena převážně na cloudové prostředí Google Cloud.

\subsubsection*{Grafické karty jako primární hardware pro inference}
\label{subsubsec:llm-computational-cost-gpu}

GPU jsou navrženy pro masivně paralelní výpočty, kde moderní karta jako NVIDIA RTX 4090 disponuje 16~384 CUDA jádry optimalizovanými pro maticové operace.
Tomu odpovídá výrazně vyšší výkon při inferenci.
Model s 7 miliardami parametrů zpracuje dotaz s 8~000 vstupními tokeny na GPU s 24 GB VRAM typicky za 15 až 60 sekund v závislosti na délce generovaného výstupu.
Výkon inference se standardně měří metrikou \textit{tokenů za sekundu} (tokens/s), která udává, kolik tokenů model vygeneruje za jednotku času.
Na moderním GPU dosahují 7B modely typicky 30 až 80 tokenů za sekundu, zatímco na CPU se tato hodnota pohybuje řádově v jednotkách tokenů za sekundu.
Dalším klíčovým faktorem je šířka pásma paměti, kterou GPU mají typicky 500 až 2~000 GB/s, zatímco u CPU je to 50 až 200 GB/s.
Protože největším úzkým hrdlem inference je přenos vah modelu z paměti do výpočetních jader, je tato vlastnost pro výkon inference zásadní.

\subsubsection*{CPU inference a její omezení}
\label{subsubsec:llm-computational-cost-cpu}

Provoz modelu na CPU je technicky možný, avšak výrazně pomalejší.
Tentýž dotaz, který GPU zpracuje za 30 sekund, může výkonnému serverovému CPU (například AMD EPYC nebo Intel Xeon) trvat 5 až 15 minut, v závislosti na velikosti modelu a délce vstupu.
Pro systém s asynchronním zpracováním, kde student nečeká na výsledek okamžitě, je takováto latence přijatelná pouze v případech velmi nízkého objemu odevzdání.
Při hromadném odevzdávání úloh před deadlinem by se fronty zpracování prodlužovaly na nepřijatelné hodnoty a výkon serveru by se výrazně snížil.

\subsubsection*{Propustnost a škálování}
\label{subsubsec:llm-computational-cost-throughput-scaling}

V prostředí školního systému je relevantní nejen latence jednotlivého dotazu, ale i propustnost systému při vyšší zátěži.
Pokud třicet studentů odevzdá řešení v krátkém časovém okně, jsou požadavky zařazeny do fronty a zpracovávány postupně.
Na GPU s dostatečnou pamětí lze požadavky zpracovávat i v dávkách (batching), čímž se zvyšuje efektivita využití hardware.
Na CPU je dávkové zpracování prakticky nerealizovatelné bez výrazného prodloužení celkové doby zpracování.

Výhodou CPU inference je však výrazně větší dostupná operační paměť (RAM), která umožňuje provozovat i modely, jež by se do VRAM grafické karty nevešly.
Existuje také hybridní přístup, kdy je část vah modelu uložena v RAM a část ve VRAM, přičemž inference probíhá částečně na CPU a částečně na GPU\@.
Tato varianta umožňuje nasadit větší modely za cenu pomalejšího zpracování.

Protože systém Kelvin zpracovává analýzu asynchronně, tedy výsledky jsou dostupné s určitým zpožděním po odevzdání, nikoliv okamžitě.
Tak je pro školní prostředí latence v řádu minut přijatelná.
Klíčové je pouze to, aby výsledky byly dostupné před tím, než vyučující přistoupí k manuálnímu hodnocení.

\endinput